% !TeX root = ../../金融经济学_熊和平_习题.tex
% 源：重拍_20260929/IMG_20260929_113635.jpg、113640.jpg（印张 62--63）

\zhchapter{风险厌恶及典型的效用函数}

\begin{tcolorbox}[summary,title=本章小结]
经济人对风险的选择由其风险态度描述，一般包括风险厌恶、风险中性和风险喜好，通常以是否愿意参与期望收益为零的公平博弈来判断：不愿参与者为风险厌恶，愿意参与者为风险喜好，无所谓者为风险中性。通常假设经济人为风险厌恶；风险厌恶对应凹效用函数，由概率论中的詹森不等式可见效用凹性与风险厌恶的关系；凸效用意味着风险喜好，线性效用意味着风险中性。马科维茨给出的风险溢价常作为风险厌恶的度量之一。

不同经济人的风险厌恶程度不同。阿罗—普瑞特用泰勒展开给出风险溢价的近似，并定义绝对风险厌恶系数与相对风险厌恶系数；确定性等价是使投资者无差异于随机财富的确定财富水平。文献中有两种确定性等价的常用定义，计算时略有差别。

两种风险厌恶系数虽来自近似，但在金融中有广泛应用。绝对风险厌恶为常数可定义 CARA 效用，相对风险厌恶为常数可定义 CRRA 效用；还可定义 HARA 及 IARA、DARA、IRRA、DRRA 等性质。CRRA 是 DRRA 的一个特例。
\end{tcolorbox}

\xhpconceptitems{
    \item 风险态度，风险厌恶，风险中性，公平博弈，詹森不等式
    \item 确定性等价，风险溢价，绝对风险厌恶系数，相对风险厌恶系数
    \item 风险容忍系数，谨慎系数，CARA，CRRA，HARA，风险厌恶比较
}

% !TeX root = ../../金融经济学_熊和平_习题.tex

\begin{problem}
\problemno{3.1} 什么是风险态度？一般包括哪几种？
\end{problem}
\begin{solution}
风险态度是经济人对随机财富的偏好，通常分为风险厌恶、风险中性和风险偏好。
\end{solution}
\begin{problem}
\problemno{3.2} 风险态度与效用函数的凹凸性之间具有什么样的关系？
\end{problem}
\begin{solution}
在边际效用为正的条件下，凹效用函数表示风险厌恶，线性效用表示风险中性，凸效用表示风险偏好。
\end{solution}
\begin{problem}
\problemno{3.3} 什么是詹森不等式？该不等式与风险厌恶有何关系？
\end{problem}
\begin{solution}
若 $u$ 为凹函数，则詹森不等式为 $u(E X)\ge E[u(X)]$。因此风险厌恶者偏好确定的期望财富，二者差额就是风险带来的效用损失。
\end{solution}
\begin{problem}
\problemno{3.4} 风险厌恶有哪几种度量方式？各种度量方法之间有何区别与联系？
\end{problem}
\begin{solution}
绝对风险厌恶系数为 $A(w)=-u''(w)/u'(w)$，以财富单位度量风险；相对风险厌恶系数为 $R(w)=-wu''(w)/u'(w)=wA(w)$，以财富比例度量风险。前者适合比较相同金额风险，后者适合比较相同比例风险。
\end{solution}
\begin{problem}
\problemno{3.5} 基于风险厌恶系数可以定义哪几种常见的效用函数？
\end{problem}
\begin{solution}
常见族包括绝对风险厌恶系数恒定的 CARA 指数效用、相对风险厌恶系数恒定的 CRRA 幂效用，以及二次效用和包含它们的 HARA 族。
\end{solution}
\begin{problem}
\problemno{3.6} 如何比较不同经济人之间的风险厌恶程度？
\end{problem}
\begin{solution}
在给定财富水平比较 $A(w)$ 或 $R(w)$；系数较大者更风险厌恶。若要作全局比较，应要求一个经济人的系数在所有相关财富水平都不小于另一个经济人。
\end{solution}

\begin{problem}
\problemno{3.7} 证明，对于 HARA 偏好：当 $\eta=0$ 时，HARA 效用函数变成了 CRRA 效用函数；当 $\gamma\to\infty$ 时，HARA 效用函数变成了 CARA 效用函数；当 $\gamma=-1$ 时，HARA 效用函数变成了二次效用函数。
\end{problem}
\begin{solution}
在 HARA 标准形式 $u(w)=\frac{(w+\eta)^{1-\gamma}}{1-\gamma}$($\gamma\neq1$)下，$\eta=0$ 时为 CRRA 幂效用；令 $\gamma\to\infty$ 并作适当标准化，极限为 $-\exp\{-aw\}$ 型 CARA 效用；$\gamma=-1$ 时指数为 $2$，得到二次效用函数。
\end{solution}
\begin{problem}
\problemno{3.8} 证明定理 3.1 中 (2) 与 (3) 等价。
\end{problem}
\begin{solution}
因 $u'(w)>0$，有 $A(w)=-u''(w)/u'(w)>0\Leftrightarrow -u''(w)>0\Leftrightarrow u''(w)<0$，所以两条件等价。
\end{solution}
\begin{problem}
\problemno{3.9} 用泰勒展开分别推导两种风险厌恶系数 ARA 和 RRA。
\end{problem}
\begin{solution}
对公平风险 $g$ 展开：$E u(w+g)=u(w)+u''(w)\operatorname{Var}(g)/2$。令 $u(w-\pi)=E u(w+g)$，得 $\pi\simeq A(w)\operatorname{Var}(g)/2$；除以财富后得到相对风险溢价 $\pi/w\simeq R(w)\operatorname{Var}(g)/(2w)$。
\end{solution}
\begin{problem}
\problemno{3.10} 已知指数效用函数 $u(w)=-\exp\{-aw\}$($a>0$)：
\begin{enumerate}[label=(\alph*)]
    \item 此函数是否表明边际效用为正且风险厌恶？
    \item 此函数是否意味着风险厌恶递减？
    \item 此效用函数是否意味相对风险厌恶为常数？
\end{enumerate}
\end{problem}
\begin{solution}
$u'(w)=ae^{-aw}>0$ 且 $u''(w)=-a^2e^{-aw}<0$，故边际效用为正且风险厌恶。$A(w)=a$ 为常数，$R(w)=aw$ 随财富线性增加，因此绝对风险厌恶恒定、相对风险厌恶不恒定。
\end{solution}
\begin{problem}
\problemno{3.11} 假设你具有 CRRA 偏好，相对风险厌恶系数为 $2$。如果你现在有 $50/50$ 的概率盈利或亏损 \num{1000} 元，同时规避风险的保险成本为 \num{500} 元，在什么财富水平下，投机活动与支付保险之间无差异？即确定性等价的财富是多少？
\end{problem}
\begin{solution}
取与 $R=2$ 等价的效用函数 $u(w)=-1/w$。无差异条件为
\[
u(w-500)=\frac12u(w-1000)+\frac12u(w+1000),
\]
整理无差异方程
\[
\frac1{w-500}=\frac{w}{w^2-1000^2}
\]
得 $w=2000$。因此确定性等价财富为 $w-500=1500$。
\end{solution}
\begin{problem}
\problemno{3.12} 经济人的初始财富为 $w$，而且他的绝对风险厌恶系数为常数 $a$，他需要承担风险 $\tilde{g}$。$\tilde{g}$ 是公平博弈，当 $\tilde{g}$ 具有如下分布时分别计算该经济人的确定性等价，并分析确定性等价与初始财富 $w$ 的关系：
\begin{enumerate}[label=(\alph*)]
    \item 取值为 $-6$ 和 $6$ 的二项分布；
    \item $[-c,c]$ 上的均匀分布；
    \item 均值为 $0$、方差为 $\sigma^2$ 的正态分布。
\end{enumerate}
\end{problem}
\begin{solution}
CARA 效用下
\[
CE=w-\frac1a\log E(e^{-a\tilde g}).
\]
二项分布 $\{-6,6\}$ 时 $CE=w-a^{-1}\log\cosh(6a)$；均匀分布 $[-c,c]$ 时 $CE=w-a^{-1}\log[\sinh(ac)/(ac)]$；正态分布时 $CE=w-a\sigma^2/2$。三种情形中风险成本均与 $w$ 无关。
\end{solution}
\begin{problem}
\problemno{3.13} 经济人的初始财富为 $w$，且绝对风险厌恶系数为 $1$，他需要承担风险 $w\tilde{g}$。$\tilde{g}$ 是公平博弈，当 $\tilde{g}$ 具有如下分布时分别计算该经济人的确定性等价，并分析确定性等价与初始财富 $w$ 的比值：
\begin{enumerate}[label=(\alph*)]
    \item 取值为 $-b$ 和 $b$ 的二项分布，其中 $0<b<1$；
    \item $[-c,c]$ 上的均匀分布，其中 $0<c<1$。
\end{enumerate}
\end{problem}
\begin{solution}
此时 $CE=w-\log E(e^{-w\tilde g})$。二项分布给 $CE=w-\log\cosh(wb)$，均匀分布给 $CE=w-\log[\sinh(wc)/(wc)]$。因此 $CE/w$ 分别取决于 $wb$、$wc$，不再是与财富无关的常数。
\end{solution}
\begin{problem}
\problemno{3.14} 经济人的初始财富为 $w$，且相对风险厌恶系数为 $2$，他需要承担风险 $\tilde{g}$，$\tilde{g}$ 取值为 $-b$ 和 $b$，$0<b<w$。有两种风险溢价的定义：
\[
    \pi_1(w)=\frac{E[u(w+\tilde{g})]}{u'(w)},\qquad \pi_2(w)=\frac{u(w)-E[u(w+\tilde{g})]}{u'(w)}.
\]
请问这两种方法计算的风险溢价相等吗？
\end{problem}
\begin{solution}
取 $u(w)=-1/w$。代入两种定义可得
\[
\pi_1=-\frac{w^3}{w^2-b^2},\qquad \pi_2=-\frac{wb^2}{w^2-b^2}.
\]
二者不相等；通常采用第二种定义，因为它对应使确定性财富下降到风险彩票期望效用的确定性金额。
\end{solution}
\begin{problem}
\problemno{3.15} 经济人的初始财富为 $w$，且是严格风险厌恶的。$\tilde{g}_1$ 和 $\tilde{g}_2$ 是两个独立的公平博弈，具有相同的分布，均为取值为 $-b$ 和 $b$ 的二项分布，其中 $0<b<w/2$。
\begin{enumerate}[label=(\alph*)]
    \item 假定他必须承担风险 $\tilde{g}_1$，则他的期望效用为 $V_1=E[u(w+\tilde{g}_1)]$。证明风险使他的情况恶化，即 $V_1=E[u(w+\tilde{g}_1)]<u(w)$。
    \item 假定他必须承担风险 $\tilde{g}=0.5(\tilde{g}_1+\tilde{g}_2)$，记在此情形下他的期望效用为 $V_2=E[u(w+\tilde{g})]$。请问 $V_2$ 是否总比 $V_1$ 高？
\end{enumerate}
\end{problem}
\begin{solution}
(a) 由严格詹森不等式及 $E\tilde g_1=0$，有 $V_1=E[u(w+\tilde g_1)]<u(w)$。(b) 平均风险 $\tilde g=(\tilde g_1+\tilde g_2)/2$ 的均值相同而分布更集中；严格凹性因此给出 $V_2>V_1$。
\end{solution}

